\Askhsh{36815}{4}
\begin{ylh}{1}
    \item 1.2 Συναρτήσεις
    \item 1.8 Συνέχεια συνάρτησης
    \item 2.4 Ρυθμός μεταβολής.
\end{ylh}
Έστω $f$ μια συνεχής συνάρτηση στο διάστημα $[-2,2]$, για την οποία ισχύει
\[ f^2(x) + x^2 = 4 \text{ για κάθε } x \in [-2,2] \]

\begin{alist}
\item Να βρείτε τις ρίζες της εξίσωσης $f(x) = 0$. \monades{06}
\item Αν η γραφική παράσταση της $f$ διέρχεται από το σημείο $A(0,2)$, τότε να βρείτε τον τύπο της $f$. \monades{09}
\item Να σχεδιάσετε τη γραφική παράσταση της $f$. \monades{04}
\item Ένα κινητό κινείται κατά μήκος της καμπύλης της $f$. Καθώς περνάει από το σημείο $B(-1, \sqrt{3})$, η τεταγμένη του $y$ αυξάνεται με ρυθμό $2$ μονάδες το δευτερόλεπτο. Να βρείτε τον ρυθμό μεταβολής της τετμημένης $x$ του κινητού τη χρονική στιγμή που περνάει από το $B$. \monades{06}
\end{alist}